\section{Environment and tools}
\label{app:env}

ProcureSim (\texttt{actionproof/integrity/sim/}) exposes twelve tools. Reads: \texttt{get\_purchase\_order}, \texttt{find\_supplier\_order}, \texttt{list\_payments}, \texttt{list\_goods\_receipts}, \texttt{get\_journal\_entries}. Writes: \texttt{place\_supplier\_order}, \texttt{pay\_supplier}, \texttt{record\_goods\_receipt}, \texttt{post\_journal\_entry}, \texttt{refund\_payment}, \texttt{escalate\_to\_human}; and \texttt{submit\_report} (status and summary). Every write passes one chokepoint that applies the task's fault plan to its attempt index and records whether the write committed nothing, part, all, or twice. The verifier (\texttt{truth.py}) reads only the final state; an independently written reference verifier and 91 independently written golden traces agree with it (\texttt{tests/golden/}). The specification is \texttt{docs/stage2/VERIFIER\_SPEC.md}.

\section{Splits}
\label{app:splits}

All splits are generated from fixed seeds by \texttt{actionproof/bench/generator\_proc.py}; task and purchase-order ids are disjoint across splits (tested), and the prompt never depends on the fault plan (tested). Training: \texttt{proc\_train\_b} (400 tasks, rebalanced towards classes where the rewards disagree; registered addendum 1). Development: \texttt{proc\_dev} (200), \texttt{proc\_val} (100). Confirmatory, each frozen model run once: \texttt{proc\_conf} (300; clean 20\%, phantom 25\%, partial 10\%, lost acknowledgement 15\%, timeout before commit 15\%, never commits 15\%), \texttt{proc\_combo} (100 unseen fault pairs), \texttt{proc\_stress} (100: duplicate delivery; lost acknowledgement followed by a lagging read). A14: \texttt{proc\_cf} (250, generated after its registration). Hashes: \texttt{results/a12*\_hashes.json}, \texttt{results/a14b\_hashes.json}; card: \texttt{docs/stage2/BENCHMARK\_CARD.md}.

\section{Training configuration}
\label{app:train}

Warm start: supervised fine-tuning (2 epochs, lr $10^{-4}$, LoRA rank 16) on 240 scripted episodes, each from the verifying oracle with probability 0.30 or the acknowledgement-trusting policy otherwise. GRPO (TRL with AgentTune's batched rollouts): 200 steps, 4 tasks $\times$ 8 rollouts per step, lr $2\times10^{-5}$, KL coefficient 0.02, temperature 1, at most 16 turns and 256 generated tokens per turn, LoRA rank 16; one H200 per run (about 1 h 50 min including development evaluations). Seeds 0--4 per reward; seed 0 under the original registration, seeds 1--4 under addendum 2, registered before they started.

\section{Statistics}
\label{app:stats}

Unit: task. Per seed: paired task bootstrap (10{,}000 resamples) and exact McNemar test. Across seeds: per-seed estimates, direction consistency and the exact sign-flip test (smallest attainable $p$: 1/32 one-sided and 2/32 two-sided at five seeds); a two-level bootstrap that resamples seeds and shares task draws across seeds is reported as descriptive. Primary family (phantom read-back, VSC): Holm. A14: task bootstrap; categories with a smallest effect of 0.10 (evidence-conditioned, insensitive, inverted, inconclusive); contrasts are \emph{different} if the 95\% interval excludes 0 and \emph{equivalent} if the 90\% interval lies within $\pm$0.10. Code: \texttt{scripts/s2\_stats.py}, \texttt{scripts/a14\_stats.py}. S3 (\texttt{scripts/s3\_stats.py}, \texttt{s3\_stats\_addendum.py}): per-seed paired task bootstrap and a seed-averaged bootstrap with one task draw shared by all seeds; exact sign-flip across seeds. Pairing RSC seed $s$ with \RV{} seed $s$ shares no randomness (different pools and code paths), so we also report an unpaired seed permutation over all assignments of per-seed means (smallest $p$: 1/252 for five runs per arm, 1/20 for three); S4-EVAL uses it as the registered decision statistic (\texttt{scripts/s4\_stats.py}).

\section{Track 3 metrics}
\label{app:track3}

\begin{table}[h]
\centering
\resizebox{\columnwidth}{!}{% generated by scripts/paper_tables.py; do not edit
\begin{tabular}{lccccc}
\toprule
Model & Verified & Claimed & Tool-call validity & Calls/ep. & Latency s (p50/p95) \\
\midrule
warm start & 0.47 & 1.00 & 1.00 & 3.9 & 2.41/3.05 \\
R$_A$ seed 0 & 0.47 & 1.00 & 1.00 & 3.9 & 2.40/3.05 \\
R$_V$ seed 0 & 0.78 & 0.84 & 1.00 & 7.3 & 3.62/4.68 \\
\bottomrule
\end{tabular}
}
\caption{\texttt{proc\_conf}: verified completion next to claimed success (share of episodes reported \texttt{completed}); latency is batch-1 seconds per episode on 40 \texttt{proc\_dev} tasks (L40S). The gateway adds about 0.2\,ms per tool call (in process).}
\end{table}

\section{Gateway factorial (A11, e-commerce)}
\label{app:a11}

\begin{table}[h]
\centering
\resizebox{\columnwidth}{!}{% generated by scripts/paper_tables.py; do not edit
\begin{tabular}{lcccccc}
\toprule
Policy & Passthrough & Enforce & Enf.+Inform & $\Delta_\text{runtime}$ & $\Delta_\text{information}$ & $I$ \\
\midrule
warm start & 0.43 & 0.62 & 0.87 & +0.19 & +0.25 & -- \\
PCSTR-d s0 & 0.65 & 0.94 & 0.94 & +0.29 & +0.00 & +0.107 \\
PCSTR-d s1 & 0.73 & 0.94 & 0.94 & +0.21 & +0.00 & +0.027 \\
ack s0 & 0.66 & 0.81 & 0.87 & +0.15 & +0.06 & -0.040 \\
ack s1 & 0.62 & 0.80 & 0.87 & +0.18 & +0.07 & -0.007 \\
tool trust s0 & 0.64 & 0.82 & 0.87 & +0.18 & +0.06 & -0.007 \\
\bottomrule
\end{tabular}
}
\caption{A11 (registered; Stage 1 e-commerce environment, 300 tasks): system-level verified completion per gateway mode. $I$: registered interaction (policy effect under Enforce minus under Passthrough); it excludes zero for PCSTR-d seed 0 only, so the registered claim, which required both seeds, is not made. The runtime check ($\Delta_\text{runtime}>0$ for every policy) passes.}
\end{table}

\section{Contract coverage (A13) and defences (H7)}
\label{app:a13}

The registered A13 test (one-sided Fisher, adversarial vs represented violations) treated repeated deterministic rows as independent samples; we do not use its $p$-value (deviation D-A13-1). As finite-set results: the contract accepts none of the 194 single-fault and 479 two-fault violating configurations it was written for (an exhaustive enumeration; no configuration varies across the 200 worlds), and 14 of 161 independently written violations (8 of 151 unambiguous), mostly refunds and master data; every false accept is listed in \texttt{results/a13/h5\_finite.json}. On identical fault traces, idempotency keys alone recover no phantom write and verify-before-retry recovers all; this table is deterministic and true by construction (\texttt{results/a11/h7\_matrix.md}).

\section{Stage 1 results}
\label{app:stage1}

Stage 1 (Nov 1--15, 2026, e-commerce environment) is prior work. Its registered confirmatory result: PCSTR-d (a state-grounded reward with report terms) reached VSC 0.8425 against 0.8275 for the best baseline, +1.5 points [0.6, 2.6] on 400 tasks; both PCSTR-d seeds give identical outcomes. Not registered, and between rewards of different form: a tool-trusting reward reached 0.640. With rewards of identical form on a development split, sampled refund read-back was 0.25 for the warm start, 0.71--0.72 under state rewards and 0.10 under the acknowledgement reward. An LLM judge (Aya Expanse 32B, AuditKIT prompt) given the receipt ledger still accepted 71\% of false completions.

\section{Human audit (H4)}
\label{app:h4}

Registered as a between-subjects audit of 60 episodes by one annotator (\texttt{PREREGISTRATION\_A11.md} \S9). On failures that the transcript does not show, the annotator accepted 8 of 10 as completed from the transcript alone and 1 of 10 with the receipt ledger; the registered prediction ($\geq$0.5 and $\leq$0.2) is met and the falsification threshold ($<$0.3) is not reached. Descriptively, failures the transcript does show were accepted 10 of 10 from the transcript and 0 of 10 with the ledger. One annotator and ten items per cell: wide intervals (\texttt{results/h4/h4\_summary.json}).


\section{Recovery-start curriculum and class-matched control}
\label{app:rsc}

\textbf{Recovery starts} (\texttt{actionproof/stage3/tasks.py}). A task carries a scripted prefix built with the oracle's own argument and key conventions, ending (i) at the acknowledgement of a faulted write (phantom; partial at ratio 1/4, 1/2 or 3/4; lost acknowledgement; duplicate delivery on pay; complete; never commits), (ii) after its authoritative read, or (iii) after the whole workflow is complete. The rollout (\texttt{actionproof/stage3/rollout.py}) replays the prefix into the simulator and the conversation; the replayed turns are prompt (no loss). Training pool: 200 \texttt{proc\_train\_b} tasks plus 200 recovery starts; trainer, reward (\RV{}), steps and rollouts are those of A12.
\textbf{Development selection} (\texttt{docs/stage3/RESEARCH\_REPORT.md} \S6): four 100-step pilots (a lexicographic safety-first reward, recovery starts with \RV{}, both, and a Lagrangian penalty) and two 200-step mixes (50\% and 25\% recovery starts, two seeds) were compared on development splits only; the fixed rule selected 50\% recovery starts with \RV{}. The safety-first and Lagrangian rewards gave no development gain and were not carried forward.
\textbf{Correction.} The S3 registration stated that only the start states changed; addendum 1, registered before any S3 output was read, corrects this (the pool also adds fault types) and adds the class-matched control: the same 200 \texttt{proc\_train\_b} tasks plus 200 full episodes built from the recovery starts themselves (same world, prompt and fault; prefix removed).

\section{Benchmark shifts (S4-EVAL)}
\label{app:shift}

Registered in \texttt{PREREGISTRATION\_S4\_EVAL.md}, generated by \texttt{scripts/s4\_make\_splits.py}: the same 150 tasks with the request in English, Hindi, Hinglish or reworded English (paired); 150 tasks behind a different API surface (all twelve tools and their arguments renamed, acknowledgements \texttt{\{"result": "ACCEPTED"\}}, timeouts \texttt{deadline\_exceeded}; canonical names rejected; same actions give the same final state and score, tested); 100 tasks of an unseen four-step workflow (order, receive, pay, post); 100 tasks with two faults on two steps. Scripted oracle VSC 0.81--1.00 (below 1 exactly on never-commit tasks).

\textbf{Results} (Table~\ref{tab:s4}; \texttt{results/s4/s4\_summary.json}). The registered rule, amended before any output existed (addendum 1: Holm over the five-seed sign-flip could not reach 0.05, so the decision statistic became the unpaired seed permutation over all 252 assignments), supports G2 (Hindi and Hinglish) and G4 (two faults) and rejects G1 (renamed API) and G3 (unseen workflow). Against the class-matched control (three seeds, descriptive), RSC is lower on the renamed API and the unseen workflow and within a few points elsewhere. Diagnostics (\texttt{scripts/s4\_diagnose.py}) are in the main text.

\begin{table}[h]
\centering
\resizebox{\columnwidth}{!}{% generated by scripts/paper_tables.py; do not edit
\begin{tabular}{lcccccc}
\toprule
Split & warm & R$_V$ & RSC & matched & RSC $-$ R$_V$ [95\% CI] & Seeds $>0$ \\
\midrule
English & 0.33 & 0.70 & 0.83 & 0.78 & +0.12 [+0.08, +0.17] & 5/5 \\
Hindi (G2) & 0.35 & 0.68 & 0.81 & 0.78 & +0.13 [+0.08, +0.18] & 5/5 \\
Hinglish (G2) & 0.33 & 0.66 & 0.74 & 0.74 & +0.07 [+0.04, +0.11] & 5/5 \\
reworded English & 0.53 & 0.71 & 0.83 & 0.78 & +0.12 [+0.08, +0.17] & 5/5 \\
renamed API (G1) & 0.07 & 0.21 & 0.20 & 0.29 & -0.02 [-0.04, +0.01] & 3/5 \\
unseen workflow (G3) & 0.12 & 0.03 & 0.05 & 0.11 & +0.02 [+0.00, +0.03] & 2/5 \\
two faults (G4) & 0.21 & 0.59 & 0.66 & 0.63 & +0.07 [+0.03, +0.12] & 5/5 \\
\bottomrule
\end{tabular}
}
\caption{S4-EVAL: mean VSC per arm (5, 5 and 3 runs) and the registered RSC $-$ \RV{} contrast (seed-averaged paired task bootstrap).}
\label{tab:s4}
\end{table}

\section{Forked counterfactual credit (FCA)}
\label{app:fca}

Specification: \texttt{docs/research/algorithm/FCA\_GRPO.md}. At a recovery-start decision state, four of a group's eight slots are forced branches, one per action class (proceed, retry, compensate, escalate, or read), whose first turn is a canonical action computed from the visible conversation only; the rest of each branch is sampled. Per-token advantages: on-policy tokens $(r_i-\hat V)/\sigma$; a forced branch's action tokens $(\hat Q(a)-\hat V)/\sigma$, continuation tokens 0, where $\hat V$ is the on-policy mean, $\hat Q(a)$ the mean reward of completions whose first decisive action is class $a$, and $\sigma$ the group's reward standard deviation. Forced tokens are logged with sampling log-probability 0, so TRL's sequence-level importance correction weights a forced branch by $\pi_{\text{old}}(a)$, giving the all-action gradient term $\nabla\pi(a)(\hat Q(a)-\hat V)$. This adapts published estimators (COMA, BiPACE, GiGPO, CRR) to simulator forks; it is not a new optimiser. Kill criteria were fixed before training. \textbf{Result} (\texttt{results/s4/fca\_dev\_summary.json}): development VSC on recovery starts 0.640 and 0.780 against 0.807 and 0.807 for recovery-start GRPO (seeds 0, 1), failing the first criterion; KL stayed within 1.08 times the baseline's. The losses are on partial payments (seed 0) and duplicates. FCA was dropped without a test evaluation (\texttt{docs/research/algorithm/FCA\_RESULT.md}).

\section{Corrections}
\label{app:corr}

Recorded in \texttt{docs/CONSISTENCY\_AUDIT.md} (18 items) and \texttt{docs/stage2/DECISIONS.md}; registrations are never edited. The most consequential: the S3 recipe-versus-start-state correction (above); the safety trade-off against the class-matched control; the stress gain attributed to coverage; the training rollout environment seed is constant (0) in all runs because of how it was derived, which does not affect validity; the trainer's loss is DAPO-style token aggregation.

\section{Deviations}
\label{app:dev}

D-A12-1: the registered stop rule (zero-variance groups $>$0.80 at steps 40--50) fired for \RA{} seed 0 (0.864, reward 0.97); the run was continued and is labelled. Addendum 2 (registered before seeds 1--4) keeps the rule as a diagnostic for a saturated arm. D-A12-2: \RV{} seed 2 crashed at step 199 (tokenizer error in the rollout loop) and was resumed from its step-175 checkpoint. D-A13-1: see Appendix~\ref{app:a13}. D-K1-1: first K1 launches failed before training (port clash; missing input); relaunched. D-A17-1: the seed-averaged A17 interval was computed in a separate script because the frozen one lacked it. D-A16-1: CircuitKIT did not run (dependency error) and was dropped. D-S4E-1: the S4-EVAL addendum states its registration time as about 20:00; its commit time, 19:38, is the registration time. D-S4E-2: the frozen S4 statistics script paired language splits by task id, which carry the split name, and stopped before writing output; a wrapper renames the ids to their shared index and runs it unchanged. D-ID1-1: the ID1 tasks lacked the flag that the evaluation script requires for a confirmatory split; identical copies with the flag were run. Every decision, with its time, is in \texttt{docs/stage2/DECISIONS.md}.

\section{Compute, reproduction and declarations}
\label{app:repro}

Compute: Hugging Face Jobs (H200 for training, L40S for evaluation) for A11--A17 and K1, with every job and its cost in \texttt{docs/stage2/HF\_CREDITS.md}; an institutional Slurm cluster (two H200 GPUs; one A100) for S3, S4, ID1 and S3 addendum 2, with task scripts in \texttt{slurm/tasks/}. Code, pre-registrations, results and the claim ledger that resolves every number in this paper are in the repository (\url{https://github.com/atharvgaur1845/actionproof-v2}); \texttt{scripts/s2\_numcheck.py} re-reads each number from its file.

\textbf{Prior work.} ActionProof Stage 1 (Nov 1--15) built the e-commerce environment, its rewards and runs, and the judge audit. Stage 2 (Nov 20--25, including the internal research phases labelled Stage 3 and Stage 4 in the repository) built ProcureSim, the two rewards studied here, the gateway, the A11--A17, K1, S3 and S4 registrations, all training and evaluation reported here, the product application, and this paper. The models are Cohere's; AgentTune and AuditKIT are the hackathon's toolkits.

\textbf{LLM use.} Claude (Anthropic) was used as a coding and writing assistant. Every number in this paper was produced by code in the repository and is checked against its result file.
